\chapter[Conclusiones]{Conclusiones}
\label{cp:conclusiones}
% Extensión estimada: 5--6 páginas, cuatro secciones y una tabla. Sin información nueva: cada cifra remite al capítulo
% en que se midió.
% Feedback del tutor (14/09/2026): al eliminarse la sección "Contribuciones" de la introducción, los aportes
% del trabajo se enuncian acá, ya respaldados por los resultados.
% Estructura tomada de la tesina de referencia (objetivos cumplidos, perspectivas futuras), con las respuestas a las
% preguntas de investigación al comienzo, como anuncia el Capítulo 6.

Este capítulo cierra el trabajo. La Sección~\ref{sec:concl_generales} responde las dos preguntas de investigación
planteadas en la Sección~\ref{sec:intro_objetivos} y reúne los hallazgos que las sostienen. La
Sección~\ref{sec:concl_objetivos} analiza en qué medida se cumplió cada objetivo, y la
Sección~\ref{sec:concl_aportes} enuncia los aportes del trabajo, cada uno con el capítulo que lo respalda. Por último,
la Sección~\ref{sec:concl_futuro} plantea las líneas de trabajo que se desprenden de las limitaciones discutidas en la
Sección~\ref{sec:disc_limitaciones}.

\section{Conclusiones generales}
\label{sec:concl_generales}
% Extensión estimada: 1.5 páginas, cinco párrafos (ficha 7.1).
% Fuentes: sec:res_clasif_discusion, sec:res_gen_discusion, sec:disc_representacion, sec:disc_texto_imagen,
% sec:disc_concordancia, sec:disc_limites. Ajustes acordados con el autor en la reestructuración del Cap. 6: PI1
% «depende del uso», no «negativa»; la comparación de DermLIP con una cabeza es entre dos formas de aprender la
% representación y no se sometió a prueba; PI2 sin tono tajante, con la verificación parcial de la descripción.

El trabajo partió de dos preguntas: si los \glspl{vlm} clasifican las lesiones de piel tan bien como los clasificadores
entrenados solo con imágenes, y si los reportes que genera un \gls{vlm} afirman el diagnóstico correcto y describen la
lesión de la imagen. Las dos se respondieron sobre la prueba de \emph{HAM10000} con una partición disjunta por lesión,
con métricas automáticas y con el cómputo de una sola GPU, y sus respuestas valen dentro de esas condiciones.

La respuesta a la primera pregunta depende de cómo se use el \gls{vlm}. Como clasificador autónomo, ninguno de los
evaluados alcanza a los supervisados. Los generativos de propósito general quedan cerca del azar, porque su respuesta
depende sobre todo del \emph{prompt}, y DermLIP sin ajuste queda entre 0.15 y 0.18 de sensibilidad macro por debajo de
los mejores supervisados, con diferencias que las pruebas estadísticas respaldan
(Capítulo~\ref{cp:res_clasif}). Como codificador de un modelo especializado, en cambio, se acerca: con una cabeza
entrenada sobre la representación congelada de DermLIP, el F1 macro llega a 0.705, frente a 0.733 de Swin-Tiny, una
diferencia que no se sometió a prueba. Esa cercanía indica que la especialización en dermatología está en la
representación, pero no cumple la motivación de la pregunta, que era obtener el diagnóstico del mismo modelo que
escribe el reporte (Sección~\ref{sec:disc_representacion}). En el melanoma, la clase de mayor prioridad clínica, la
diferencia la marca la regla de decisión más que el modelo: el umbral fijado en validación lleva la detección de
Swin-Tiny de tres a cuatro de cada cinco casos.

La respuesta a la segunda pregunta es que el reporte afirma el diagnóstico correcto en la medida en que lo hace el
clasificador que lo ancla. Sin un diagnóstico, ningún generador afirma de manera confiable la clase de la imagen: el
que más la nombra lo hace enumerando dos clases por reporte, y en ninguno más de uno de cada veinte reportes nombra
únicamente la clase verdadera. Con la cascada, en cambio, el generador transcribe el diagnóstico que recibe, de modo que la consistencia
diagnóstica de Mini-VLM, 0.849, coincide con la exactitud de Swin-Tiny, 0.847, y con un diagnóstico equivocado el
reporte lo afirma en lugar de corregirlo (Capítulo~\ref{cp:res_gen}). El ruteo por concordancia eleva la consistencia
a 0.870 y, entre los diagnósticos que afirma, reduce la proporción de errores de 0.153 a 0.085
(Sección~\ref{sec:disc_concordancia}). En cuanto a la descripción, la imagen cambia lo que el reporte describe de cada
lesión, porque sin ella la cobertura de entidades clínicas de Mini-VLM cae a la mitad, pero la descripción también se
adapta al diagnóstico recibido, y que cada rasgo descrito esté en la imagen queda verificado solo en parte.

Las dos respuestas se encadenan. Como un \gls{vlm} generativo no decide la clase por sí mismo, el diagnóstico lo pone
un clasificador; y como el generador lo transcribe, la fidelidad diagnóstica del reporte queda fijada por ese
clasificador y no por el generador (Sección~\ref{sec:disc_texto_imagen}). En consecuencia, el valor del generador
está en la redacción y en la descripción de cada lesión, y el de la concordancia con el \gls{knn}, en señalar los casos
en que el diagnóstico es menos confiable. Lo que el sistema no resuelve son los errores que comparten los dos
predictores, que se afirman con la misma plantilla que un acierto y que incluyen casi uno de cada cuatro melanomas. Por
eso el \emph{pipeline} puede ordenar la atención del especialista, pero no reemplazarla.

Por último, el trabajo muestra que las decisiones de evaluación pesan tanto como la elección del modelo. La partición
por imagen sobreestima la sensibilidad macro en unas 0.07, una cifra del orden de la distancia con los mejores
resultados de la literatura, y la variación entre ejecuciones supera una décima en las clases minoritarias, mientras
que Swin-Tiny, ConvNeXt-Tiny y ViT-B/16 no se distinguen entre sí (Sección~\ref{sec:disc_limites}). Esa observación,
junto con el costo medido, hace compatible la restricción de una GPU de consumo con la fidelidad diagnóstica del
\emph{pipeline}: Mini-VLM ocupa como máximo 8.6~GB y escribe un reporte en unos 16 segundos, y la fidelidad depende de
un clasificador pequeño.

\section{Cumplimiento de los objetivos}
\label{sec:concl_objetivos}
% Extensión estimada: 1.5 páginas, una tabla y tres párrafos (ficha 7.2).
% Fuentes: sec:intro_objetivos (OE1--OE3, tab:intro_objetivos_preguntas), balances de los Caps. 4 y 5,
% sec:disc_limitaciones. El grado de cumplimiento se declara con lo que no se hizo, sin atenuarlo.

El objetivo general era desarrollar y evaluar un sistema, basado en modelos preentrenados y adaptados, capaz de
asistir en el diagnóstico de lesiones de piel y en la generación de sus reportes, y analizar cuán bien clasifica, si
el reporte afirma el diagnóstico correcto y si puede ejecutarse con recursos convencionales. El sistema se construyó
como un \emph{pipeline} de tres componentes, un clasificador, un \gls{knn} sobre la representación de DermLIP y un
generador, y cada una de esas tres preguntas recibió una respuesta medida. La Tabla~\ref{tab:concl_objetivos} resume,
para cada objetivo específico, lo que se hizo, el resultado principal y los capítulos en que se presenta.

\begin{table}[H]
    \centering
    \caption{Objetivos específicos: lo que se hizo, el resultado principal y dónde se presenta.}
    \label{tab:concl_objetivos}
    \small
    \setlength{\tabcolsep}{4pt}
    \begin{tabularx}{\textwidth}{@{}>{\raggedright\arraybackslash}p{2.2cm}
        >{\raggedright\arraybackslash\hspace{0pt}}X >{\raggedright\arraybackslash\hspace{0pt}}X
        >{\raggedright\arraybackslash}p{1.7cm}@{}}
        \toprule
        \textbf{Objetivo} & \textbf{Lo que se hizo} & \textbf{Resultado principal} & \textbf{Capítulos} \\
        \midrule
        OE1. Clasificadores & Cuatro supervisados, dos \glspl{vlm} generativos y DermLIP, evaluados con un mismo
            protocolo y con pruebas estadísticas, los supervisados con tres ejecuciones; umbral de melanoma fijado en validación
            & Swin-Tiny, con F1 macro de $0.733 \pm 0.026$; ningún \gls{vlm} sin una cabeza entrenada alcanza a los
            supervisados; el umbral detecta cuatro de cada cinco melanomas
            & \ref{cp:metodos}, \ref{cp:res_clasif} \\
        \filasep
        OE2. \gls{vlm} para el reporte & Mini-VLM ajustado con LoRA en el codificador visual sobre \emph{SkinCaRe}, comparado con
            LLaVA-Med y SkinVL, sin diagnóstico, en cascada con el clasificador y con el ruteo por concordancia
            & El reporte afirma el diagnóstico que recibe; la consistencia diagnóstica es de 0.849 con la cascada y
            de 0.870 con el ruteo
            & \ref{cp:metodos}, \ref{cp:res_gen} \\
        \filasep
        OE3. Escenarios y costo & Prueba externa en \emph{HIBA}, perturbaciones, calibración, casos dudosos
            separados por el ruteo, métricas del texto y costo de cada componente
            & Pérdida de cerca del 40\,\% de la sensibilidad macro en \emph{HIBA}; riesgo selectivo de 0.085 entre
            los afirmados; Mini-VLM entra en 8.6~GB y escribe un reporte en unos 16 segundos
            & \ref{cp:res_clasif}, \ref{cp:res_gen}, \ref{cp:discusion} \\
        \bottomrule
    \end{tabularx}
\end{table}

El primer objetivo se cumplió en los términos en que se planteó. Los clasificadores de las tres familias se
diseñaron, se implementaron y se evaluaron con la misma partición por lesión y las mismas métricas, y la detección de
melanoma se analizó por separado, con un umbral elegido en validación y no en la prueba. Esa comparación es la que
permite responder la primera pregunta de investigación.

El segundo objetivo se cumplió con una modificación que los resultados justifican. El \gls{vlm} ajustado produce el
reporte e integra la imagen en la descripción, pero no produce por sí mismo un diagnóstico preliminar confiable: sin
un diagnóstico, no afirma de manera confiable la clase de la imagen (Sección~\ref{sec:res_gen_sin_anclaje}). Por eso,
en el sistema final el diagnóstico lo aporta el clasificador y el generador lo transcribe en el reporte. El sistema,
como un todo, produce el diagnóstico preliminar y el reporte, pero el diagnóstico no proviene del \gls{vlm}.

El tercer objetivo se cumplió en parte. Se evaluaron los escenarios que el objetivo nombra: la otra población, con
\emph{HIBA}, los casos en que el sistema no tiene seguridad, con el ruteo, la calidad del texto, con las métricas
automáticas, y el costo computacional, en una GPU de consumo. Quedaron sin evaluar, en cambio, dos aspectos. El
\emph{pipeline} completo no se aplicó a \emph{HIBA}, de modo que la fidelidad del reporte en otra población se infiere
de la caída del clasificador y no se midió. Además, la utilidad práctica se midió solo por el costo: ningún
especialista leyó los reportes, y el efecto del sistema sobre el flujo de trabajo clínico, que el objetivo general
menciona, no se evaluó (Sección~\ref{sec:disc_limitaciones}).

\section{Aportes del trabajo}
\label{sec:concl_aportes}
\pendiente{Se enuncian los aportes, cada uno con remisión al capítulo que lo sustenta: una comparación
sistemática bajo un mismo protocolo entre clasificadores supervisados, VLMs médicos generales y un modelo
visión--lenguaje dermatológico; el diagnóstico del colapso de los VLMs generales en clasificación
dermatológica y la evidencia de que el factor determinante es la especialización de dominio y no la
modalidad; un ensamble en cascada consciente de la incertidumbre que mejora la fidelidad diagnóstica de los
reportes, replicado con dos generadores; la auditoría de fuga de información entre particiones de HAM10000 y
los resultados corregidos con una partición disjunta por lesión; y el análisis de viabilidad sobre hardware
de consumo.}

\section{Trabajo futuro}
\label{sec:concl_futuro}
% Extensión estimada: 1 página, una oración de entrada, una lista de cinco líneas y un párrafo de cierre (ficha 7.4).
% Fuentes: sec:disc_limitaciones (alcance, ancla de una ejecución, decisiones sin evaluar), sec:disc_concordancia,
% sec:disc_texto_imagen, sec:disc_limites. Sin futurología (guía §3): cada línea dice qué falta evaluar, no qué
% beneficio se espera.

Las líneas de trabajo futuro se desprenden de las limitaciones de la Sección~\ref{sec:disc_limitaciones}, y cada una
consiste en evaluar algo que este trabajo dejó sin medir:
\begin{itemize}
    \item \textbf{La evaluación por especialistas.} La continuación más directa es que dermatólogos lean los reportes
    junto con sus imágenes, para verificar que los rasgos descritos estén en la lesión y juzgar si el reporte sería
    aceptable en la práctica. Esa evaluación, junto con un conjunto de reportes de referencia redactados por
    especialistas para las imágenes dermatoscópicas, permitiría pasar de la fidelidad diagnóstica a la validez
    clínica.
    \item \textbf{Una regla de decisión orientada al costo clínico.} El ancla del \emph{pipeline} usa la clase más
    probable de Swin-Tiny, no el umbral de melanoma, y el \emph{pipeline} se eligió por una métrica que trata igual
    todos los errores (Sección~\ref{sec:disc_concordancia}). Queda por evaluar cómo cambian la consistencia por clase y
    la proporción de casos derivados si el ancla usa ese umbral, o si el ruteo pondera los errores según su costo.
    \item \textbf{Otras señales de confianza.} Los errores del clasificador y del \gls{knn} coinciden unas tres veces
    más de lo que coincidirían por azar, posiblemente porque los dos aprenden de las mismas imágenes. Queda por
    evaluar si la probabilidad del clasificador, que está bien calibrada, o un conjunto de casos de referencia
    distinto del de entrenamiento detectan una parte de esos errores compartidos.
    \item \textbf{La validación externa del \emph{pipeline}.} En \emph{HIBA} se evaluó solo el clasificador. Aplicar
    el \emph{pipeline} completo a otra población, y a otras modalidades como la fotografía clínica, permitiría medir la
    fidelidad del reporte fuera de \emph{HAM10000} en lugar de inferirla.
    \item \textbf{Otros anclajes y generadores.} La relación entre la exactitud del clasificador y la fidelidad del
    reporte se caracterizó con tres fuentes y una sola ejecución de Swin-Tiny. Un barrido de clasificadores y de
    ejecuciones la describiría con más detalle. Del mismo modo, quedan por evaluar los generadores que razonan antes de
    diagnosticar y los de mayor tamaño, junto con una medida de qué parte de la descripción proviene de la imagen y
    qué parte del diagnóstico recibido (Sección~\ref{sec:disc_texto_imagen}).
\end{itemize}

Todas estas líneas parten de la misma conclusión: en un sistema que combina un clasificador y un generador, la
fidelidad diagnóstica del reporte depende del clasificador y de la señal que decide cuánto afirmar. Por eso las
mejoras que este trabajo deja planteadas apuntan a esos dos componentes y a su evaluación con especialistas, más que al
tamaño del generador.

